\section{Why: two refuted hypotheses and one real cause}
\label{sec:why}

\subsection{Divergence is not the cause}

The standard explanation is lockstep execution: lanes run in groups of 32, most edges fail
their guard early, and one survivor drags 31 idle lanes through the remaining pipeline. We
instrumented lane survival directly:

\begin{center}
\begin{tabular}{lrrr}
\toprule
Model & Work items & Survivors & Survive \\
\midrule
Viking N=15 & $8\,355\,840$ & $8\,355\,840$ & \textbf{100.0\%} \\
Fischer N=10 & $3\,317\,950$ & $2\,372\,190$ & 71.5\% \\
\bottomrule
\end{tabular}
\end{center}

On Viking there is \emph{no} early-exit divergence whatsoever. Stream compaction or wavefront
scheduling --- the standard remedy, and the idea we were about to implement --- would recover
$0\%$ there and at most $28\%$ on Fischer. Moreover the only DBM-free filter, the integer
guard, is already applied host-side (Fig.~\ref{fig:kernel}), so a cheap pre-pass has nothing
to reject. \emph{We report this specifically so that others do not implement it.}

\subsection{Arithmetic intensity is not (yet) the cause}

Each successor moves ${\approx}1152$\,B for ${\approx}2600$ integer operations:
$2.3$\,ops/byte against a roofline ridge point of ${\approx}10$, so the kernel is memory-bound
\emph{in character}. But the achieved rates are nowhere near that bound:

\begin{center}
\begin{tabular}{lrrr}
\toprule
& Rate & \% compute peak & \% bandwidth \\
\midrule
GPU expand (1400\,ms) & 5.97\,M succ/s & \textbf{0.5\%} & 2.5\% \\
CPU expand (491\,ms) & 17.0\,M succ/s & ${\approx}20\%$ (of CPU) & 7\% \\
\bottomrule
\end{tabular}
\end{center}

The bandwidth ceiling for this kernel is ${\approx}237$\,M succ/s, roughly $40\times$ above
what the GPU achieves. The GPU is limited by \emph{utilisation}, not by a physical roof: the
CPU runs at ${\approx}20\%$ of its peak while the GPU runs at $0.5\%$ of its own.

\subsection{The working-set dilemma}
\label{sec:dilemma}

Because Floyd--Warshall's $k$ loop is strictly sequential, there is almost no parallelism
\emph{inside} one zone, so every thread must own a whole zone. That creates a dilemma with no
good side, and we measured both horns (Table~\ref{tab:reps}).

\begin{table}[t]
\centering
\caption{Expansion time by representation. R3 cuts per-lane footprint $24\times$ and is the
slowest of the three.}
\label{tab:reps}
\begin{tabular}{llrr}
\toprule
& Representation & Viking N=15 & Fischer N=10 \\
\midrule
R1 & private dense (spills to device memory) & \textbf{1425\,ms} & \textbf{183\,ms} \\
R2 & threadgroup-resident dense & 1587\,ms & 277\,ms \\
R2$'$ & \quad + bank-conflict padding & 1575\,ms & 297\,ms \\
R3 & row-distributed across SIMD lanes & 2536\,ms & 351\,ms \\
\bottomrule
\end{tabular}
\end{table}

Keeping the zone \textbf{private} means 2304\,B per thread, far beyond the ${\approx}128$--$256$
registers a thread receives, so it spills to thread-local \emph{device} memory and all
${\approx}2600$ DBM operations per successor run at device-memory latency. Keeping it
\textbf{on-chip} costs $32 \times 576\,\text{B} = 18$\,KB of Apple's 32\,KB threadgroup
budget, leaving roughly one resident threadgroup per core: occupancy collapses and there is
nothing left with which to hide latency. R1 wins because sheer thread count buys more latency
hiding than on-chip residency buys in access speed.

\subsection{No representation escapes it}
\label{sec:cooperation}

R3 was designed to break the dilemma and does so on its own terms --- $\mathit{dim}$ ints per
lane, comfortably register-resident --- yet it is the slowest. The reason is a counting
argument, and it generalises to \emph{any} cooperative mapping.

One zone per \emph{thread} keeps $32$ independent zones in flight per SIMD group. Any scheme
that cooperates \emph{within} a zone --- R2's shared tile or R3's distributed rows --- processes
$1$ zone per SIMD group. Cooperation must therefore make a single zone more than $32\times$
faster merely to break even. The parallelism available inside one zone is only
$\mathit{dim}$-wide: at $\mathit{dim}=12$ that is 12 of 32 lanes busy, an immediate
$2.7\times$ loss, with a ceiling of $12\times$ before paying for \texttt{simd\_shuffle}
traffic and the serial $k$ loop. \textbf{Twelve-fold best case against a thirty-two-fold
deficit.} Both cooperative attempts lost, by different mechanisms, exactly as predicted.

This is the deepest reason a ``GPU-friendly representation'' does not exist here: the
workload's parallelism lives \emph{across} zones and is already fully exploited by the
simplest possible mapping. The dense matrix inherited from \uppaal{} is not what holds the
GPU back.

\subsection{Unified-memory coherency}

Finally, a hybrid in which the CPU expands and the GPU merges --- attractive because the two
phases favour different hardware --- is defeated by cache coherency. Identical merge work
costs 258\,ms when the zones were written by the GPU and 571\,ms when written by the CPU,
because the GPU must snoop them out of CPU caches. The moment the CPU produces the zones, the
GPU is the wrong consumer.
